%==============================================================================
%  Correction to "Post-Quantum Extension of the Conditional CRO Trilemma"
%  IEEE Access 14 (2026) 125130, DOI 10.1109/ACCESS.2026.3722742
%
%  Note: this draft uses the standard article class so that it compiles
%  anywhere. For submission, paste the body into the official IEEE Access
%  template (ieeeaccess.cls), keeping the section structure unchanged.
%==============================================================================
\documentclass[10pt,a4paper]{article}
\usepackage[margin=2.2cm]{geometry}
\usepackage{amsmath,amssymb,amsthm,mathtools}
\usepackage{microtype}
\usepackage[hidelinks]{hyperref}

\newtheorem*{lemmaSeven}{Lemma 7}
\newtheorem*{theoremTwo}{Theorem 2}
\newtheorem*{assumptionFourP}{Assumption 4$'$}
\newtheorem*{assumptionSix}{Assumption 6}
\theoremstyle{remark}
\newtheorem*{remark}{Remark}

\newcommand{\Confq}{\mathsf{Conf}^q}
\newcommand{\Relq}{\mathsf{Rel}^q}
\newcommand{\Oppq}{\mathsf{Opp}^q}
\newcommand{\negl}{\mathsf{negl}}
\newcommand{\cP}{\mathcal{P}}
\newcommand{\IMS}{I(M;\Sigma\mid C)}
\newcommand{\Iacc}{I_{\mathrm{acc}}(M;\Sigma\mid C)}

\title{\textbf{Correction to ``Post-Quantum Extension of the Conditional CRO
Trilemma: Impossibility Bounds on Confidentiality, Reliability, and
Opposability Under Quantum Institutional Extraction Limits''}}
\author{Thierry Emmanuel Minka Mi Nguidjo\"{\i}$^{1,2}$,
        Flavien Serge Mani Onana$^{1}$,
        and Thomas Djotio Ndi\'e$^{1}$\\[4pt]
\small $^{1}$Laboratory of Mathematical Engineering and Information Systems,
National Advanced School of Engineering,\\
\small University of Yaound\'e I, Yaound\'e, Cameroon\\
\small $^{2}$Supreme State Audit Office, Yaound\'e, Cameroon\\[2pt]
\small Corresponding author: Thierry Emmanuel Minka Mi Nguidjo\"{\i}
(minkathierry@gmail.com)}
\date{}

\begin{document}
\maketitle

\noindent In the above article~\cite{pq}, the proof of Theorem~2
(Post-Quantum Conditional CRO Trilemma) inherits the errors corrected
in~\cite{cor1} for the classical Conditional CRO Trilemma~\cite{cro}:
Step~4 uses the classical statement of Lemma~4 of~\cite{cro}, whose
inequality has the wrong direction; Step~3 and Assumption~4 use
Assumption~2 in the wrong direction; and the constant $C(\beta_q,\alpha_q)$
of Assumption~5 originates from the incorrectly solved optimization step
of~\cite{cro}. In addition, inequality~(68) in the proof of Lemma~6 does
not hold as written, and Remark~17 overstates the range of validity of
Theorem~2. These errors were identified by the authors during the
machine-checked formalization of the results in Lean~4. This correction,
which parallels~\cite{cor1}:
(i)~withdraws Lemma~6 and the upper bound~(72) of Lemma~7, and restates the
lower bound~(71) with the Holevo information;
(ii)~states the verdict-efficiency condition explicitly (Assumption~$4'$);
(iii)~states explicitly the leakage regime in which Theorem~2 holds
(Assumption~6) and gives a complete proof with explicit constants; and
(iv)~corrects Remark~17 and clarifies the reading of Section~VI-C/D.
The quantum institutional model (Assumptions~1--3), the institutional
parameters $(\Delta_q,\alpha_q,\beta_q,\eta_q)$, the quantum Fano step
(81)--(85), and the form of the main bound (80) are preserved within the
stated regime. The authors apologize for any inconvenience caused.

%------------------------------------------------------------------------------
\section{Lemmas 6 and 7 (Section V-A)}
%------------------------------------------------------------------------------

\paragraph{Lemma 6.} Inequality~(68), with the expectation taken over
\emph{independent} uniform draws of $m_0,m_1$, does not hold in general:
for two messages encoded in orthogonal (classical, perfectly
distinguishable) states, $I_{\mathrm{acc}}=1$ bit, whereas the right-hand
side of~(68) equals $1/2$. Lemma~6 is therefore withdrawn. The derivation
(78)--(79) of the upper bound~(72) of Lemma~7 depends on Lemma~6 and is
also withdrawn; note that~(78) is moreover inconsistent with the linear
form~(70). Neither Lemma~6 nor~(72) is used in the corrected proof below.

\paragraph{Lemma 7.} The lower bound~(71) follows from the quantum Pinsker
inequality and quantum regularity (Assumption~3), with one adjustment: the
``simplified form''~(74) replaces the Holevo information
$\chi(\mathcal E_{m,c})$ of Assumption~3 by the accessible information,
whereas the Holevo bound gives $I_{\mathrm{acc}}\le\chi(\mathcal E_{m,c})$,
so that this substitution goes in the wrong direction. Stated with the
Holevo information, the lower bound holds and is retained:

\begin{lemmaSeven}[Quantum confidentiality guarantee; replaces Lemma 7]
For any $(\beta_q,\eta_q)$-quantum-regular protocol $\cP$ (Assumption~3),
\begin{equation*}\tag{71$'$}\label{eq:l7}
\Confq(\cP)\;\ge\;1-\frac12\sqrt{2\beta_q\,\chi_{\max}+2\eta_q\,D_{\max}}\,,
\end{equation*}
where $\chi_{\max}=\max_{c}\chi(\mathcal E_{m,c})$ and
$D_{\max}=\max_{c}D(\bar\rho^{\,c}_\Sigma\,\|\,\mathbb I/d_\Sigma)$; for the
NIST PQC schemes of Remark~13, $D_{\max}=\negl(\lambda)$.
\end{lemmaSeven}

\begin{proof}
For all $m_0,m_1,c$, the quantum Pinsker inequality~(73) and Assumption~3 give
$\|\rho^{m_0,c}_\Sigma-\rho^{m_1,c}_\Sigma\|_1^2\le
2\beta_q\chi(\mathcal E_{m,c})+2\eta_qD(\bar\rho^{\,c}_\Sigma\|\mathbb I/d_\Sigma)
\le 2\beta_q\chi_{\max}+2\eta_qD_{\max}$.
Taking the maximum over $(m_0,m_1,c)$ in~(76) gives~\eqref{eq:l7}.
\end{proof}

\noindent This is the quantum counterpart of Lemma~4 as corrected
in~\cite{cor1}. Remark~16 is modified accordingly: only the lower
bound~\eqref{eq:l7} is used.

%------------------------------------------------------------------------------
\section{Verdict efficiency (Assumption 4 and Step 3)}
%------------------------------------------------------------------------------

Assumption~2 gives $\log_2 k_\lambda\le\alpha_q H(C)$, an \emph{upper}
bound. Hence neither the last step of~(58) nor the second inequality
of~(90) follows from Assumptions~2 and~4. The lower bound on $H(V\mid C)$
is the efficiency property of institutional verdict functions, which is now
stated explicitly.

\begin{assumptionFourP}[Verdict non-degeneracy and efficiency; replaces Assumption 4]
The institutional verdict function uses the available context efficiently:
\begin{equation*}\tag{57$'$}\label{eq:eff}
H(V\mid C)\;\ge\;\theta_q\,\alpha_q H(C)\;\eqqcolon\;\gamma_q,
\qquad \gamma_q\ge 1,
\end{equation*}
for a non-degeneracy constant $\theta_q\in(0,1]$.
\end{assumptionFourP}

\noindent With $\theta_q=1/2$ and the context entropies of Table~2
($H(C)\ge 9$ bits, $\alpha_q\ge 0.8$), one has $\gamma_q\ge 3.6$.
Claim~1 remains valid through its second argument
($\alpha_qH(C)\ge H(V\mid C)\ge 1$).

%------------------------------------------------------------------------------
\section{Corrected statement and proof of Theorem 2 (Section V-B)}
%------------------------------------------------------------------------------

The dominant-term analysis of Step~6 assumes $u=x/\lambda=O(\lambda^{-1})$
(``since $x\le C_{\mathrm{proof}}=O(1)$''), while Remark~17 asserts that the
bound holds for all $I\in[0,\lambda]$, including $I=\Theta(\lambda)$. The
latter assertion is incorrect: for $u$ bounded away from $0$, the factor
$(\lambda u+\Delta_q)/(\alpha_qH(C))$ in~(96) grows linearly with $\lambda$.
As in~\cite{cor1}, the leakage regime is therefore stated explicitly.

\begin{assumptionSix}[Confidentiality-preserving regime]
There exists a constant $\kappa\ge 0$, independent of $\lambda$, such that
\begin{equation*}\tag{59$'$}\label{eq:regime}
\chi_{\max}=\max_{c}\chi(\mathcal E_{m,c})\;\le\;\frac{\kappa}{\lambda}.
\end{equation*}
\end{assumptionSix}

\noindent Since $\IMS$ is the average of $\chi(\mathcal E_{m,c})$ over the
contexts and $\Iacc\le\IMS$ by the Holevo bound, Assumption~6 implies
$\Iacc\le\IMS\le\kappa/\lambda$. By Lemma~7, it is the regime in which
quantum-regular protocols guarantee confidentiality:
$\Confq\ge1-\tfrac12\sqrt{2\beta_q\kappa/\lambda+2\eta_qD_{\max}}$, which tends
to $1$ when $D_{\max}=\negl(\lambda)$.

\begin{theoremTwo}[Post-Quantum Conditional CRO Trilemma, corrected]
Under Assumptions~1--3, Assumption~$4'$ \eqref{eq:eff} and
Assumption~6 \eqref{eq:regime}, for any post-quantum evidence protocol
$\cP$ with security parameter $\lambda\ge64$, message entropy
$H(M)=\lambda$, $M$ independent of $C$, and quantum institutional
parameters $(\Delta_q,\alpha_q,\beta_q,\eta_q)$:
\begin{equation*}\tag{80$'$}\label{eq:main}
\Confq(\cP)\cdot\Relq(\cP)\cdot\Oppq(\cP)
\;\le\;\frac{1}{\gamma_q}\left(\frac{\Delta_q}{\lambda}
+\frac{\kappa(\Delta_q+1)}{\lambda^{2}}+\frac{\kappa^{2}}{\lambda^{3}}\right)
\;\le\;\frac{\Delta_q+\eta_q H(C)}{\lambda}+O(\lambda^{-2})+\negl(\lambda).
\end{equation*}
\end{theoremTwo}

\begin{proof}
Write $x=\IMS$; by Assumption~6, $x\le\kappa/\lambda$.

\emph{Step 1 (unchanged).} By the quantum Fano inequality, (81)--(85) give
$\Iacc\ge\Relq\cdot\lambda-1$ for every $\Relq\in[0,1]$. Since
$\Iacc\le x$,
\begin{equation*}\tag{85$'$}
\Relq(\cP)\le\frac{x+1}{\lambda}.
\end{equation*}

\emph{Step 2.} By Definition~23, Assumption~1~(87) and
Assumption~$4'$,
\begin{equation*}\tag{89$'$}
\Oppq(\cP)=\frac{I(V;\Sigma\mid C)}{H(V\mid C)}\le\frac{x+\Delta_q}{\gamma_q}.
\end{equation*}

\emph{Step 3 (product bound).} Since $0\le\Confq(\cP)\le1$ and all factors
are nonnegative,
$\Confq\cdot\Relq\cdot\Oppq\le\frac{x+1}{\lambda}\cdot\frac{x+\Delta_q}{\gamma_q}$.
The right-hand side is nondecreasing in $x$; since $x\le\kappa/\lambda$,
\[
\Confq\cdot\Relq\cdot\Oppq\le\frac1{\gamma_q}
\Bigl(\frac1\lambda+\frac{\kappa}{\lambda^2}\Bigr)
\Bigl(\Delta_q+\frac{\kappa}{\lambda}\Bigr)
=\frac1{\gamma_q}\Bigl(\frac{\Delta_q}{\lambda}
+\frac{\kappa(\Delta_q+1)}{\lambda^2}+\frac{\kappa^2}{\lambda^3}\Bigr).
\]
The last inequality in~\eqref{eq:main} follows from $\gamma_q\ge1$ and
$\eta_qH(C)\ge0$.
\end{proof}

\noindent Steps~3--7 of the original proof, equations (90)--(111), and
Assumption~5 with the constant $C(\beta_q,\alpha_q)$ are replaced by the
above; Assumption~5 is no longer needed. The quantum interpretability loss
$\eta_q$ enters the confidentiality guarantee (Lemma~7), the channel
bounds (Lemma~10), and the regime classification of
Section~VI-A, which are unaffected.

\begin{remark}[Replaces Remark 17]
Assumption~6 is essential. For the classical diagonal encoding of
Example~3 (SPHINCS+, $\beta_q=\beta$, $\eta_q=0$), the protocol that reveals
$M$ with probability $1/2$ and otherwise outputs a uniformly random proof
satisfies Assumptions~1--3 with $\beta_q\le1.04$ for all $\lambda\ge64$ and
has $\Confq=1/2$, $\Relq=1/2+2^{-\lambda-1}$ and, with $V$ the first bit of
$\Sigma$, $\Oppq=1$, for every $\lambda$; here
$\IMS=\lambda/2-1+O(\lambda2^{-\lambda})$ lies outside
regime~\eqref{eq:regime}. The bound therefore does not hold uniformly for
$I=\Theta(\lambda)$.
\end{remark}

%------------------------------------------------------------------------------
\section{Clarification of Section VI-C/D and Remark 10}
%------------------------------------------------------------------------------

The implied capacities of Tables~8--10, the Dilithium-3 case study
(Section~VI-D) and item~2 of Remark~10 are obtained by inverting~(80) on
observed or projected CRO products $P_{\mathrm{obs}}$. Under Assumption~6,
Step~1 gives $\Relq(\cP)\le(1+\kappa/\lambda)/\lambda$ in the sense of
Definition~22, so these products, computed on the operational CLO scale,
are not instances of~\eqref{eq:main}. The implied values of $\Delta_q$
should therefore be read as an empirical calibration on the CLO scale,
consistent with the bounded-rationality range of Section~IV-D, and not as
consequences of Theorem~2. The operational regimes of Section~VI-A, the
sensitivity analysis of Section~VI-B and the design recommendations of
Section~VIII are unaffected.

\begin{thebibliography}{9}
\bibitem{pq}
T.~E. Minka Mi Nguidjo\"{\i}, F.~S. Mani Onana, and T.~Djotio Ndi\'e,
``Post-quantum extension of the conditional CRO trilemma: Impossibility
bounds on confidentiality, reliability, and opposability under quantum
institutional extraction limits,'' \emph{IEEE Access}, vol.~14,
pp.~125130--125155, 2026, doi: 10.1109/ACCESS.2026.3722742.
\bibitem{cro}
T.~E. Minka Mi Nguidjo\"{\i}, F.~S. Mani Onana, and T.~Djotio Ndi\'e,
``The conditional CRO trilemma: Impossibility bounds on confidentiality,
reliability, and opposability under institutional extraction limits,''
\emph{IEEE Access}, vol.~14, pp.~33263--33274, 2026,
doi: 10.1109/ACCESS.2026.3669065.
\bibitem{cor1}
T.~E. Minka Mi Nguidjo\"{\i}, F.~S. Mani Onana, and T.~Djotio Ndi\'e,
``Correction to `The conditional CRO trilemma: Impossibility bounds on
confidentiality, reliability, and opposability under institutional
extraction limits','' \emph{IEEE Access}, submitted, 2026.
\end{thebibliography}

\end{document}
